% !TeX root = ../example.tex
\begin{frame}[fragile,label=demo-notes]{正文与备注分工}
  \begin{columns}[T,onlytextwidth]
    \begin{column}{.48\textwidth}
      \textbf{页面源码}
      \begin{verbatim}
\begin{itemize}
  \item 核心要点
\end{itemize}
\note{解释、过渡和讲述提示。}
      \end{verbatim}
      备注与页面写在同一份 TeX 中。
    \end{column}
    \begin{column}{.48\textwidth}
      \begin{description}[PDF 备注屏]
        \item[默认 PDF] \texttt{shownotes=false}，只显示幻灯片。
        \item[PDF 备注屏] \texttt{shownotes=true}，备注位于右侧。
        \item[PPTX 备注] 从 TeX 导出逐页备注，供演讲者视图使用。
      \end{description}
    \end{column}
  \end{columns}
\end{frame}
\note{讲者提示：正文保留核心信息，解释、过渡和讲述提示写在 note 中。shownotes 是 documentclass 选项，默认关闭；开启后 PDF 是左侧幻灯片加右侧备注的双屏布局，需要支持这种布局的阅读器。PPTX 始终从同一份 TeX 备注导出，PDF 是否显示备注不会改变 PPTX 备注来源。逐步显示页面可用 note<1>、note<2> 为各阶段分别写备注。}
